\begin{proposition}[Positive Euler number for rational normalization]\hypertarget{tag-02WT}{}\label{night:P4.21}\label{stmt:P4.21}
\StatementRecord{night:P4.21}{proved}{02WT}{\SeriesRef{cdp:st:A.5.3}, \SeriesRef{night:P4.12}, \SeriesRef{night:P4.13}, \SeriesRef{night:P4.20}, \SeriesRef{night:P4.6}, \SeriesRef{night:P4R.1}}

Let $f:Y\to B$ be a proper surjective holomorphic map with connected
fibres from a smooth compact complex threefold with $b_1(Y)=b_2(Y)=0$
to a smooth compact curve, with smooth general fibre a complex two-torus.
Let $mS$ be a fibre with irreducible singular reduced support, and let
$\nu:T\to S$ be its smooth rational normalization. Then $e(S)>0$.
\end{proposition}
\begin{proof}

\begin{enumerate}[label={$\langle1\rangle$\arabic*.},leftmargin=*]
\item Let $D$ be the conductor divisor on $T$, with its scheme
multiplicities, and put $C=\operatorname{Sing}(S)_{\mathrm{red}}$.
Propositions~\ref{cdp:st:A.5.3} and~\ref{night:P4.12} give
$\omega_T(D)\simeq\nu^*\omega_S$, a torsion line bundle.
The minimal resolution in the first proposition is the identity of the
smooth surface $T$. Pullback to a common point-blowup resolution of
$T$ and $\mathbb P^2$ injects $\operatorname{Pic}(T)$ into a free
abelian Picard group; see \cite[II.3, pp.~17--18; II.12, p.~24]{Beauville1978}.
Thus $D\sim-K_T$. Contract $T$ to a minimal rational surface $T_0$,
which is $\mathbb P^2$ or a Hirzebruch surface $\mathbb F_n$ with
$n\ne1$ \cite[V.10, pp.~79--80]{Beauville1978}, and push $D$ forward
to $D_0$. For each reversed point blowup $q$ with exceptional curve $E$,
the canonical formula $K_{\mathrm{new}}=q^*K_{\mathrm{old}}+E$
gives $D_{\mathrm{new}}=q^*D_{\mathrm{old}}-E$.
Its exceptional coefficient is $\mu-1$, where $\mu$ is the
conductor multiplicity at the centre. Effectivity therefore puts
every centre on the conductor.

\item On $\mathbb P^2$, the divisor $D_0$ is a plane cubic. An
irreducible cubic is smooth elliptic, rational nodal, or
rational cuspidal. A reducible cubic is a smooth conic
and a line, three lines, a double line and another line,
or a triple line. The conic and line meet at two
transverse points or at one tangency; three distinct
lines form a triangle or are concurrent. Thus either
$D_0$ is reduced nodal and has fundamental group
$\mathbb Z$ or $\mathbb Z^2$, or its reduced support is
simply connected and all normalized components are
rational.

\item On $\mathbb F_n$, write $h$ for the tautological class
and $F$ for a ruling fibre. The formulas
$h^2=n$, $h\cdot F=1$, $F^2=0$, and
$K_{\mathbb F_n}=-2h+(n-2)F$ are
\cite[III.18, pp.~46--47; IV.1, pp.~53--54]{Beauville1978}.
Thus $D_0\cdot F=2$, and its horizontal part is a double
section, two distinct reduced sections, or a reduced
irreducible bisection. A double section together with
vertical fibres is a rational tree. For the other cases
write
\[
 D_0=H_1+H_2+\sum_j a_jF_j
 \quad\hbox{or}\quad
 D_0=H+\sum_j a_jF_j.
\]
In the first case adjunction on either section gives
\[
 H_1\cdot H_2+\sum_j a_j=2.
\]
The support is therefore a rational tree or a single
rational cycle. A cycle occurs only with all coefficients
one and ordinary double intersection points. This
includes the cases of two section-intersection points,
one section-intersection point and one vertical fibre,
or two vertical fibres. Coincident intersections or a
vertical fibre through the section-intersection point
instead give trees.

\item For a bisection, adjunction gives
\[
 p_a(H)=1-\sum_j a_j.
\]
There is either no vertical component and $H$ has
arithmetic genus one, or one coefficient-one vertical
fibre and a smooth rational bisection. In the former
case $H$ is smooth elliptic, rational nodal, or rational
cuspidal. In the latter it meets the fibre at two distinct
points or at one tangency. Again the only nontrivial
fundamental groups are $\mathbb Z$ and $\mathbb Z^2$,
and those cases have reduced nodal conductor.

\item We follow these alternatives through the blowups. At a
point of conductor multiplicity at least two, the new
conductor contains the exceptional $\mathbb P^1$.
Its contraction is the quotient which collapses that
connected simply connected subcomplex to the old
point. Van Kampen therefore makes the induced map on
fundamental groups an isomorphism. At a point of
multiplicity one the conductor is smooth and reduced,
and its strict transform is isomorphic to it. Thus
$\pi_1(D_{\mathrm{red}})$ is unchanged throughout.
If the initial group was nontrivial, the conductor was
reduced nodal. A smooth-point blowup adds no conductor
component, and a node blowup inserts one component of
coefficient one. It remains reduced nodal.

\item In this latter case the generic transverse conductor
exponent of every normalization branch is one. The
exponent is its intrinsic conductor exponent plus the
intersection multiplicities with the other branches.
A singular branch has intrinsic exponent at least two.
Exponent one therefore means two smooth branches
meeting transversely. Consequently $S$ has ordinary double curves outside finitely
many points. Proposition~\ref{night:P4.20} makes $S$ normal crossing
everywhere. Write $t$ for its number of ordinary triple points.
On each normal-crossing chart the composite $T\to S\subset Y$
is the inclusion of a smooth coordinate plane, hence is immersive.
The exact sequence
\[
0\longrightarrow T_T\longrightarrow\nu^*(T_Y|_S)
 \longrightarrow L\longrightarrow0
\]
has line quotient. Since $H^2(Y,\mathbb Q)=H^4(Y,\mathbb Q)=0$,
the pulled-back first and second Chern classes vanish. Thus
$c_1(L)=K_T$ and $e(T)=K_T^2$. Rationality gives
$\chi(\OO_T)=1$ \cite[III.20--21, pp.~48--50]{Beauville1978};
Noether's formula therefore gives $e(T)=6$.
The reduced nodal anticanonical conductor is either smooth elliptic
or a cycle of rational curves. Adjunction on a rational component
gives exactly two node preimages. Its number of components therefore
equals its number of nodes, and its Euler number equals that number;
the smooth elliptic case has Euler number zero and no nodes. Every triple point
of $S$ has three node preimages on $D$, so $e(D)=3t$.
Over the smooth part of $C$ the conductor has two points per fibre,
whereas a triple point has three. Hence
\[
e(D)=2e(C\setminus\{\text{triple points}\})+3t
     =2e(C)+t,
\]
and $e(C)=t$. Additivity gives $e(S)=6-2t$.
Proposition~\ref{night:P4.6}, applied to this reduced divisor $S$,
gives $e(S)=t$. Therefore $t=2$ and $e(S)=2$.

\item In the remaining case $D_{\mathrm{red}}$ is connected
and simply connected, and all its normalized components
are rational. Triangulate the analytic pair $(T,D_{\mathrm{red}})$ and the finite
map $D_{\mathrm{red}}\to C$. The inclusion is a cofibration, so
its pushout has the fundamental group of the corresponding mapping
cylinder. Van Kampen gives
\[
 \pi_1(S)=\pi_1(T)*_{\pi_1(D_{\mathrm{red}})}\pi_1(C)
          =\pi_1(C).
\]
Every component of $C$ has rational normalization,
because it is the finite nonconstant image of a
rational component of $D_{\mathrm{red}}$. Proposition~\ref{night:P4R.1} makes
$\pi_1(C)$ virtually abelian. For a connected curve
with rational normalized components, this is the free
group $F_b$ of its point-identification graph, so
$b\le1$. If $c\ge1$ is its number of components, then
\[
 e(C)=c+1-b\ge c\ge1.
\]
Proposition~\ref{night:P4.13} gives $e(S_{\mathrm{reg}})\ge0$. Additivity proves
$e(S)=e(S_{\mathrm{reg}})+e(C)>0$.

\end{enumerate}
\end{proof}